% !TeX root = ../教材题库_全部条目_转写工作稿.tex
% Source: XJ-B1-C05-TH-P166-Y0848
% 第5章 / 5.2 任意角的三角函数 / 教材 p.166、167 / PDF 第173、174页
% OCR未校对：公式与文字均可能有误；未完成数学TeX转写；禁止直接发给学生
% 原题及相邻原解截图（核对用）：../题目原图/XJ-B1-C05-TH-P166-Y0848-01.webp, ../题目原图/XJ-B1-C05-TH-P166-Y0848-02.webp
% 本文件只在转写完成后作为题干使用；OCR 条目尚不是可靠题干。
\begin{bankitem}{p.166 思考（含紧邻说明）}{第5章 / 5.2 任意角的三角函数 / 教材 p.166、167 / PDF 第173、174页}{XJ-B1-C05-TH-P166-Y0848}{OCR未校对：公式与文字均可能有误；未完成数学TeX转写；禁止直接发给学生}
\begin{ocrdraft}
[教材p.166 / PDF第173页]
\par
？思考如何将正弦函数值、余弦函数值、正切函数值用几何方法直观表示？
\par
在三角函数的定义式中，正弦sina=x与余弦灬a一一的分母都是厂，其中厂
\par
十3，2．为了简单起见，我们可以取7一1，即让点P(x，回在单位圆上，此时
\par
sina=J，，cosa=“而03，都可以作线段来表示．具体作图方法如下：
\par
如图5．2一6，设单位圆的圆心为平面直角坐标系的原点0，角a的终边与单位
\par
圆交于点P，过点P作轴的垂线，垂足为D.
\par
蚱声生
\par
（2）
\par
（3）
\par
（4）
\par
图5．2一6
\par
由三角函数的定义可知，三角函数值sina=y，灬a一有正负之分，若仅用线
\par
段DP,OD来分别表示它们还不够，需引人方向．为此，我们引人有向线段的概念．
\par
规定了方向（即规定了起点和终点）的线段称为有向线段．例如，将DP看作有
\par
向线段，其中D为起点，P为终点．当DP的方向与轴的正方向相同时，其数量
\par
取正实数值3，；当DP的方向与3，轴的负方向相同时，其数量取负实数值3，；当DP
\par
[教材p.167 / PDF第174页]
\par
的长度为0时（即角a的终边落在轴上），其数量取零．于是，在所有的情况下都
\par
有DP==y==s1na\ensuremath{\cdot}
\par
将OD看作有向线段，其中0为起点，D为终点．同理可得OD一工一co艹．
\par
由单位圆与角a的终边的交点P作出的有向线段DP，它的方向和长度分别代表
\par
了sina的符号和绝对值，DP代表的实数就是角仅的正弦，故有向线段DP称为角
\par
的正弦线．同理，我们将有向线段OD称为角仅的余弦线．
\par
由正切函数的定义可知，只要在角a的终边上取一点T使它的横坐标
\par
为1，则纵坐标就等于tan如图5．2一6，过点A（1，（))作单位圆的切线工一1.当
\par
角的终边落在第一、四象限或工轴的正半轴时，该切线与终边相交于点T（1，）．
\par
N\_AT
\par
于是，有tan一
\par
而OA一1，则有tana=AT．
\par
当角a的终边落在第二、三象限或轴的负半轴时，该切线与终边的反向延长
\par
线也相交于点T(1，）．类似地，有tana=AT.
\par
因此，只要tan存在，则上述与单位圆相切的有向线段，AT（A为起点，T为终
\par
点）代表的实数就是tan，故，AT称为角的正切线．
\par
正弦线、余弦线、正切线统称三角函数线．
\end{ocrdraft}
\end{bankitem}
